\chapter{Diskussion, Limitationen, Perspektiven und Fazit}\label{diskussion-limitationen-perspektiven-und-fazit}

\section{Diskussion}\label{diskussion}

In der Analyse ist es gelungen, auf der Grundlage eines sehr großen, drei Zeitungen umfassenden Korpus und eines etwa zwei Jahrzehnte umspannenden Untersuchungszeitraums Frames und Framings einzelner Extremismusvarianten differenziert herauszuarbeiten. Mithilfe der vorgeschlagenen Methode konnten, einzelne FE mitsamt der sie evozierenden sprachlichen Zeichen datengeleitet erschlossen werden. Die Grundlage dafür waren distributionelle Strukturen der sprachlichen Oberfläche, die in korpuslinguistischer Theoriebildung schon lange als bedeutsam erkannt worden sind und durch ihre systematische Auswertung und Inbezugsetzung genutzt wurden, um Frame-semantische Einheiten zu rekonstruieren.

Als besonderes Potenzial ist der induktive und explorative Charakter der Methode sowie die Skalierbarkeit auf sehr große Datenmengen bzw. Untersuchungszeiträume hervorzuheben. Nur mithilfe computergestützter Verfahren kann es m.\,E. gelingen, dem von Busse für eine linguistische Epistemologie angesetzten Ziel einer Rekonstruktion sämtlichen verstehensrelevanten Wissens und seiner sprachlichen Realisierungen nahezukommen. Gerade wenn man berücksichtigt, dass dieses Wissen in ständigem Wandel begriffen und kontextspezifisch zu modellieren ist, erscheint seine Rekonstruktion alleine mithilfe manuell-qualitativer Ansätze geradezu aussichtslos. Unter Anwendung korpuslinguistischer Verfahren konnten semantisch reichhaltige Modelle erschlossen werden, die eine Grundlage für die analytische Auseinandersetzung darstellten -- nicht als ‚objektive` Ergebnisse, sondern im Sinne einer „methodisch abgesicherten und unterstützten, skrupulösen, empirisch verfahrenden ‚Hermeneutik`\,`` \autocite[156]{busseDiskursSpracheGesellschaftliches2013}. Der explorative Charakter der Methode ergibt sich neben der \emph{corpus-driven} Fundierung auch durch die erzeugten Visualisierungen. Mit der diagrammatischen Grundfigur des Netzwerkgraphen eröffnet sich eine Möglichkeit, die komplexen Relationen einzelner FE -- auf eine auch ästhetisch ansprechende Weise -- darzustellen.

Punktuell wurde das datengeleitete Vorgehen um theoretisch hergeleitete bzw. auf Diskurswissen gestützte Elemente komplementiert. Dies betrifft zum einen die Definition von orientierenden wesentlichen Merkmalen des Extremismuskonzeptes (\chapref{inhaltliche-mindestanforderungen-an-einen-extremismus-frame}), die es ermöglicht haben, einzelne Sprachgebrauchsmuster überhaupt als einschlägig für den Extremismusdiskurs zu erkennen; zum anderen die Selektion einzelner \emph{seed words}, um einen Punkt im Vektorraum zu berechnen, mit dessen Hilfe sich eine grobe thematische Einschlägigkeit einzelner Wort-Cluster bewerten lässt. Dass die Selektion der \emph{seed words} die erzeugten Frame-Strukturen nicht übermäßig beeinflusst hat, zeigt sich insbesondere an der Netzwerk-Community ‚Staat, FdGO \& Gesellschaft`. Während man z.\,B. das Zu-Tage-Treten der Community ‚Politik` auf das verwendete \emph{seed word} \emph{Npd} zurückführen könnte, ist das FdGO-Netzwerk-Cluster durch keines der ausgewählten \emph{seed words}, die eher typische Merkmale oder Akteure des Extremismus widerspiegeln, prädeterminiert. Zugleich stiftet die FdGO aber als Antithese zum Extremismusbegriff \emph{ex negativo} seinen definitorischen Kern. Dass dieser inhaltlich für einen Extremismus-Frame zentrale Frame-Anteil \emph{corpus-driven} emergent wurde, stellt einen schönen Beleg für das Funktionieren der Methode dar.

Die Analyse hat insbesondere die Mesostruktur der erzeugten Frames im Sinne von Extremismusvarianten-Frames herausgearbeitet und dabei makrostrukturelle Kontextualisierungen im Sinne thematischer Teil-Frames berücksichtigt. Die Mikrostruktur im Sinne einzelner FE oder LE wurde insofern angesprochen, als immer wieder wesentliche, durch einzelne oder einige wenige FE konstituierte Anteile der Varianten-Frames wie z.\,B. der Hanau- und der George-Floyd-Frame im Zeitraum 2020--2021 oder der NSU-Frame im Zeitraum 2011--2013 thematisiert wurden. Auch die Illustration von Relationen zwischen den FE anhand einzelner Konkordanzen stellt ein solches Element dar. Hier wäre es aber natürlich möglich, noch weitaus ausführlichere Detailanalysen -- seien sie qualitativer oder \emph{corpus-based} quantitativer Natur -- anzuschließen und z.\,B. aus einer semasiologischen Perspektive einzelnen Ausdrücken in ihrem sich wandelnden Perspektivierungspotenzial nachzugehen. So könnten etwa zentrale Begriffe wie (u.\,a.) \emph{Radikalismus}, \emph{Extremismus}, \emph{Totalitarismus} und \emph{Fundamentalismus} in Hinblick auf die mit ihnen assoziierten Wissensbestände in den Blick genommen werden. Dies würde einen neuen methodischen Zugang zu einer in den Sozialwissenschaften mit regem Interesse verfolgten Fragestellung (siehe \chapref{begriffliche-verwirrungen-und-kritik-am-extremismus-konzept}–\ref{weitere-arbeiten-zum-extremismusdiskurs-und--begriff}) eröffnen. In einer onomasiologischen Perspektive ließe sich auch genauer nachzuvollziehen, was mit Frame-Teilen passiert, die aus dem medialen Fokus geraten. Was genau passiert etwa, wenn im Zeitraum 2002--2010 das Gewaltpotenzial des Rechtsextremismus ausgeblendet wird? Wie werden die einzelnen Rechtsextremen Akteure im Anschluss geframet? Werden die mit ihnen assoziierten gewaltbezogenen FE bzw. LE dann vom Islamismus übernommen oder finden sich dort bei genauerer Betrachtung andere Ausdrücke? In qualitativen Analysen, die jedoch „quantitativ informiert`` \autocite{bubenhoferQuantitativInformierteQualitative2013} sein sollten, können Veränderungen des Framings außerdem an einzelnen Texten nachvollzogen werden. Dabei können dann insbesondere solche Aspekte fokussiert werden, die sich nicht (nur) auf der lexikalischen Ebene vollziehen -- etwa die Analyse von Argumentationsmustern oder rhetorischen Fragen.

Gleichwohl war es auch aus der eingenommenen Vogelperspektive möglich, wesentliche Umstrukturierungen der Extremismus-Semantik herauszuarbeiten, die etwa das Gefahrenpotenzial einzelner Extremismusvarianten, einzelne Akteure und ihre Einstellungen oder spezifische Praktiken in den Vordergrund rückten oder auch über lange Zeiträume ausblendeten. Durch die Auswahl von Leitmedien als Datengrundlage ist es nicht überraschend, dass die Veränderungen im Framing in weiten Teilen den zeitgeschichtlichen Verlauf widerspiegelt, in dem sich einzelne in der Berichterstattung aufgegriffene und so diskursiv konstituierte Ereignisse herausbilden. Mit dem Fokus auf der Ermittlung von FE und Frame-Relationen zwischen diesen FE wurden zwar wesentliche Teile von Busses Frame-Konzept modelliert und ein Beitrag zur Lösung des zentralen Problems der Granularität und Exhaustivität von Frame-Beschreibungen geleistet -- es konnten jedoch nicht alle von Busse vorgesehenen Modell-Ebenen erfasst werden. Dazu zählen etwa Hierarchien bzw. Vererbung von FE, die Annotation von weiteren Arten der Frame-Relationen, wie z.\,B. Constraints und die Identifikation von Default-Werten. Auf der Grundlage der vorliegenden Daten wäre es jedoch möglich, diese Desiderata anzugehen und etwa zu versuchen, FE zu ermitteln, die im Sinne eines Constraints nie gemeinsam auftreten. Auch zur Identifikation von Default-Werten, also solchen Teilen des Wissens, die in einzelnen Äußerungen nicht explizit gemacht werden müssen, um mitverstanden zu werden, ist ein wesentlicher Schritt mit den jeweils über große Subkorpora berechneten Frames geleistet. Selbstverständlich werden die zahlreichen FE nicht alle zugleich in einzelnen Texten realisiert -- um jedoch einen genaueren Einblick zu bekommen, wie sich Instanziierungsmuster darstellen, wäre es für folgende Arbeiten von Interesse, die Verteilung der FE in einzelnen Texten des Korpus genauer zu beleuchten. Auf weitere methodische Desiderata wird außerdem im folgenden Kapitel eingegangen.

Ein weiteres Ziel der Arbeit bestand darin, auszuloten, welchen Beitrag quantitative, \emph{corpus-driven} Methoden bei der Herausarbeitung eines komplexen epistemischen Modells wie dem semantischen Frame-Modell Busses leisten können. Dass es mit ihrer Hilfe möglich war, mit der vorliegenden diskursanalytischen Literatur kongruente Strukturen des diskursiven Wissens zu beschreiben, ohne zuvor Kategorienschemata zu entwickeln, erproben und an Einzeltexten oder (Trainings\mbox{-})Datensätzen zu annotieren, zeigt den Erfolg dieses Ansatzes und verdeutlicht, dass der häufig geäußerte Einwand, dass die ‚wahre` Bedeutung ja zwische den Zeilen liege, wo der Computer natürlich nicht lesen könne, durchaus fraglich ist.

Ein Ansinnen der vorliegenden Arbeit bestand außerdem darin, den Erfolg bei der Modellierung nicht einfach als im Rahmen einer Blackbox generierten Output ‚dankbar hinzunehmen`, sondern linguistisch zu reflektieren. Die wesentlichen verwendeten korpuslinguistischen Methoden (Word Embeddings und Kollokationen), die sich durch eine systematische Berücksichtigung und Kombination der syntagmatischen und paradigmatischen Verteilung der lexikalischen Einheiten im Korpus auszeichnen, wurden daher in die distributionelle Semantik und Korpuspragmatik eingebettet und in Hinblick auf eine Frame-semantische Modellierung operationalisiert.

Es wurde bei der Methodenentwicklung mit Kollokationen und Word Embeddings auf sehr verbreitete Algorithmen gesetzt, die in ihrer Funktionsweise vergleichsweise gut nachvollziehbar sind. Die Berechnung sprachbasierter Embeddings stellt dabei eine in der Computerlinguistik sehr etablierte Methode zur Sprachmodellierung dar, die jedoch in korpuslinguistischen Ansätzen kaum Anwendung findet. Wie auch Scharloth et al. \autocite*{scharlothWuchernRhizomeLinguistische2013} und Bubenhofer et al. \autocite*{bubenhoferDigitaleKulturlinguistikDigitalitaet2022} hervorheben, ist es jedoch von großem Interesse für die Korpuslinguistik, zum einen nicht den Anschluss an technologische Entwicklungen zu verlieren, sie zum anderen aber auch in ihrer Funktionsweise zu reflektieren und ihre Potenziale und Grenzen für linguistische Fragestellungen auszuloten. Gerade vor dem Hintergrund der neuesten technologischen Entwicklungen im Bereich der dialogfähigen LLM ist dieses Desiderat zu unterstreichen. Die vorliegende Arbeit hat versucht, verständlich zu machen, dass und warum relativ komplexe Bedeutungszusammenhänge aus der Verteilung sprachlicher Zeichen algorithmisch rekonstruierbar sind.

Während das Ziel der Arbeit nicht darin bestand, eine ‚objektive` Methode zu entwickeln, die ohne Interpretation auskommt, sondern einen quantitativ-hermeneutischen Zugang zum Diskurs zu wählen, ist nicht von der Hand zu weisen, dass der Erfolg der Erschließung der umfangreichen Frame-Strukturen vom Welt- bzw. Diskurswissen der Forschungsperson abhängig ist. Keine der ermittelten Strukturen spricht für sich, sondern muss -- ohne dabei zu viele Vorannahmen in die Daten zu projizieren und mithin den induktiven Ansatz auszuhebeln -- gedeutet werden. Die generierten Frame-Visualisierungen sind darüber hinaus sehr umfangreich und komplex, sodass die hermeneutische Auseinandersetzung durchaus Zeit benötigt. Dies ist für sich genommen kein Nachteil der Methode, denn auch Diskurse sind komplexe Gegenstände, die sich nicht ohne Weiteres auf wenige relevante Wissenselemente bzw. sie evozierende sprachliche Einheiten herunterbrechen lassen. Gerade für die wissenschaftliche diskursanalytische Auseinandersetzung ist die Komplexität daher m.\,E. gegenstandsadäquat. Nichtsdestoweniger könnte es vielversprechend sein, Möglichkeiten der Komplexitätsreduktion zu erproben, insbesondere um die Methode für andere Kontexte wie z.\,B. Didaktik im Bereich Deutsch als Fremd- und Zweitsprache fruchtbar zu machen. Dort besteht seit einiger Zeit die Forderung, Lernende „zum Verstehen deutschsprachiger Diskurse und zur Partizipation an ihnen zu befähigen, d. h. sie mit kulturbezogenem Wissen und kulturbezogenen Kompetenzen aus{[}zu{]}statten`` \autocite[184]{altmayerLandeskundeAlsKulturwissenschaft2006}. Eine offene Frage ist jedoch, wer die umfangreichen dafür vorgesehenen „rekonstruktiv-qualitativ{[}en{]}`` \autocite[191]{altmayerLandeskundeAlsKulturwissenschaft2006} Diskursanalysen zu den zahllosen sich im Wandel befindlichen aktuellen Diskursen und diskursiven Mustern durchführen sollte. Hier stellen quantitative, in Teilen automatisierte Verfahren m.\,E. einen Lösungsansatz dar.

Äußerst spannend wäre es außerdem, Synergieeffekte der hier vorgeschlagenen Methode bzw. der Analyseresultate mit \emph{FrameNet} auszuloten. Ich habe in \chapref{warum-frames} die Auffassung vertreten, dass \emph{FrameNet}-Frames für diskursanalytische und korpuspragmatische Forschungsvorhaben sowie komplexe epistemologische Modellierungen zu sehr an valenzgrammatischen und kontextabstrakten lexikographischen Ideen ausgerichtet sind. Nichtsdestoweniger wäre es nun interessant, die Gegenprobe zu machen und zu überprüfen, wie sich die hier ermittelten Strukturen zu \emph{FrameNet}-Frames verhalten: Stimmen Teile der datengeleitet erschlossenen Frames mit \emph{FrameNet}-Frames überein? Gibt es umgekehrt Teile der jeweiligen Frame-Netzwerke, die nicht erfasst sind und einander ergänzen könnten? Die Beantwortung dieser Fragen könnte nicht nur die Anbindung an aktuelle Frame-semantische Forschung stärken, sondern auch einen Einstieg zur besseren Interpretierbarkeit der hier vorgeschlagenen Frames bieten bzw. eine Möglichkeit der Komplexitätsreduktion im Sinne des zuvor formulierten Desiderats darstellen.

\section{Methodische Limitationen und Perspektiven}\label{methodische-limitationen-und-perspektiven}

Wie im vorherigen Kapitel beschrieben, wurde im Rahmen der entwickelten Methode insbesondere auf etablierte und vergleichsweise gut nachvollziehbare Algorithmen gesetzt. Aus computerlinguistischer Perspektive stellt sich jedoch die Frage, ob insbesondere mit dem verwendeten Word-Embedding-Algorithmus von Mikolov et al. \autocite*{mikolovEfficientEstimationWord2013} nicht der \emph{state of the art} verfehlt wurde, liegen doch mit den Sprachmodellen der Transformer-Generation wie BERT \autocite{devlinBERTPretrainingDeep2019} und -- noch neuer -- ChatGPT \autocite{openaiGPT4TechnicalReport2023} avanciertere und in einigen Anwendungen überlegene Algorithmen vor. Mit der Möglichkeit, kontext- bzw. kotextspezifische Token-Embeddings zu generieren, d.\,h. nicht mehr alle formseitig identischen Token gleich zu behandeln, können sie u.\,a. Homonymieprobleme lösen und stellen so eine durchaus vielversprechende Weiterentwicklung für linguistische Fragestellung dar.\footnote{Auch im Kontext der Analyse in der vorliegenden Arbeit kam es vereinzelt zu Problemen durch Type-basierte Embeddings, etwa wenn der Ausdruck \emph{Republikaner} im Cluster \cluster{99\_Republikaner{\breakbar}Demokraten} sowohl auf die US-amerikanische als auch die deutsche Partei verweist.} Ein bedeutender Nachteil dieser Modelle besteht jedoch darin, dass sie ein Vielfaches der Ressourcen im Vergleich zu den Embedding-Algorithmen der \emph{Word2Vec}-Generation beanspruchen -- sowohl in Hinblick auf den Umfang der Trainingsdaten als auch auf die verwendete Computer-Hardware. Sie können nicht auf eigenen Datensätzen trainiert werden, sodass auf vortrainierte und häufig in Hinblick auf die Datenbasis intransparente Modelle zurückgegriffen werden muss, die bestenfalls noch auf eigenen Daten nachtrainiert werden können. Dies stellt einen gewichtigen Nachteil für einen sprachgebrauchs- und kontextsensiblen bzw. korpuspragmatischen Ansatz dar. Des Weiteren ist es keinesfalls so, dass sie \emph{per se} performanter als Type-basierte Embeddings sind, sondern in einzelnen Aufgabenstellungen durchaus schlechter abschneiden \autocites[siehe etwa][]{brunnerBERTNotBERT2020, ehrmanntrautTypeTokenbasedWord2021, lenciComparativeEvaluationAnalysis2022}. Besonders interessant ist in dieser Hinsicht, dass sich zwei Ansätze beim Umgang mit Word Embeddings unterscheiden lassen. Einerseits der in der Computerlinguistik häufig verfolgte Ansatz, Embeddings als Wortrepräsentationen in einzelnen \emph{tasks} wie \emph{sentiment analysis}, Autorschaftsattribution oder auch Frame-Induktion einzusetzen. Dabei stellen die Embeddings nur eine notwendige Transformation dar, um Wörter für den Computer auswertbar zu machen, werden jedoch nicht selbst in ihrem Distributionsverhalten in den Blick genommen. Andererseits findet sich ein in den Digital Humanities oder auch in der Korpuspragmatik bzw. in dieser Arbeit praktizierter Ansatz, in dem Word Embeddings als „\emph{abstractions of semantic systems} by describing word meaning as a set of relation of a focus term to its semantic neighbours`` eingesetzt werden.\footnote{Während die Autor*innen dabei in erster Linie lexikalisch-semantische Ontologien wie \emph{GermaNet} oder den \emph{Duden} vor Augen haben, die nicht deckungsgleich mit korpuspragmatischen Sprachmodellierungsansätzen sind, entsprechen sie diesen doch deutlich eher als die erstgenannten Verfahren.} Ehrmanntraut et al. kommen zu folgendem Ergebnis:

\begin{quote}
The most important result from our experiments is that a widespread assumption in NLP and Computational Humanities is not true: a context-sensitive embedding like BERT is not automatically better for all purposes. Static embeddings like fastText are at least on par if not better if word embeddings are used as abstractions of semantic systems. \autocite[28]{ehrmanntrautTypeTokenbasedWord2021}
\end{quote}

Vor diesem Hintergrund erscheint es nicht nur zu rechtfertigen, sondern sogar durchaus empfehlenswert, sich in korpuspragmatisch-diskursanalytischen Ansätzen auf klassische Word Embeddings zu stützen. Auch dem Plädoyer von Egbert et al. (2020: 39), „appropriate \emph{and} minimally sufficient statistical methods`` anstatt übermäßig komplexe und schwer nachvollziehbare statistische Modelle zu nutzen, wird dadurch entsprochen.

Auch wenn die verwendeten Algorithmen deutlich geringere Datenmengen erfordern als Sprachmodelle der neueren und neuesten Generation, funktionieren auch sie nur bzw. deutlich besser mit großen Korpora. Es ist insofern ein wichtiges Desiderat, zu überprüfen, inwiefern in kleinen oder mittelgroßen Korpora, wie z.\,B. dem Parlamentsdebatten enthaltenden \emph{GermaParl}-Korpus \autocite{blaetteGermaParlCorpusPlenary2020}, semantische Frames mit der vorgeschlagenen Methode erschlossen werden können. Ein vielversprechender Ansatz könnte darin bestehen, ein vortrainiertes Word-Embedding-Modell zu nutzen und nur die \emph{count-based} Kollokationen auf dem kleineren Korpus zu berechnen, da diese bei kleineren Korpora verlässlichere Ergebnisse liefern. So würden zumindest die Frame-Relationen zwischen den vortrainierten FE kontextspezifisch modelliert werden. Auch das Nachtrainieren eines auf einem größeren Korpus berechneten Sprachmodells könnte erprobt werden.

Des Weiteren wäre es äußerst interessant, einzelne Diskurse gewissermaßen durch die Brille anderer Diskurse zu betrachten. Also z.\,B. Embeddings auf einem Korpus zu berechnen, ihre Äquivalente in einem anderen Korpus zu bestimmen (z.\,B. mit TWEC) und auf dieser Basis Frames zu modellieren. So könnten auch Konzepte, die in anderen Korpora gar nicht als solche bestehen, in den Blick genommen werden -- etwa ob es bereits vor der Entstehung des Extremismuskonzeptes in den Sozialwissenschaften in den 1960er Jahren semantisch äquivalente Frames bzw. ähnliche Sprachgebrauchsmuster gab.

Bezüglich der Implementierung wäre es, wie in \chapref{interpretation-der-knoten} bereits angesprochen, eine wichtige Optimierung, andere Keyword-Statistiken als die von \emph{polmineR} berechnete \emph{Log-Likelihood} zu erproben, da diese Wörter ausblendet, die ausschließlich im fokussierten Korpus, nicht aber im Referenzkorpus vorkommen, was insbesondere Neologismen und seltene Namen nicht erfasst. Sehr vielversprechend erscheint es mir weiterhin, die Erkenntnisse von Kozlowski et al. \autocite*{kozlowskiGeometryCultureAnalyzing2019} zur Messbarkeit der semantischen Multidimensionalität von Konzepten einzubeziehen. Ließe sich für jedes Wort bzw. FE berechnen, wie es auf semantischen Dimensionen wie ‚gut -- schlecht`, ‚gefährlich -- ungefährlich`, ‚links -- rechts`, ‚extremistisch -- gemäßigt` oder ‚historisch -- zeitgenössisch` verortet ist, wäre dies hochgradig informativ für die Interpretation der Cluster als FE und ein Verständnis ihres Framing- bzw. Perspektivierungspotenzials.

Des Weiteren ist aufgefallen, dass einige Cluster sehr groß und kaum interpretierbar sind bzw. häufig nur einen semantisch homogenen Kern haben, zu den Rändern hin aber vage werden. Hier könnte ausgelotet werden, ob die Cluster im Nachhinein um Ausreißer bereinigt werden können, etwa indem die enthaltenen Ausdrücke hinsichtlich auffällig großer Cosinus-Distanzen zum Centroid überprüft werden. Ein weiteres Desiderat in dieser Hinsicht besteht darin, paradigmatische Ähnlichkeit, z.\,B. von Wörtern in einem Cluster, besser nachvollziehbar zu machen. Diese Ähnlichkeit liegt in der Ähnlichkeit der einzelnen Kotexte der Wörter begründet, ist jedoch mitunter schwer durch Sichtung einzelner Belegstellen zu ergründen, gerade wenn es nicht einige Dutzend, sondern einige hundert oder tausend Treffer für die jeweiligen Ausdrücke gibt. Hier wäre es wünschenswert, einzelne illustrative Belegstellen für die paradigmatische Ähnlichkeit geclusterter Ausdrücke automatisch identifizieren zu können. Sprachoberflächliche Indikatoren könnten das Vorkommen geteilter Kollokate in den Sätzen oder die Nähe der jeweiligen Word Embeddings zu Embeddings der Satzebene (\emph{sentence embeddings}) sein.

Wie in \chapref{kenndaten-zum-korpus-und-korpusaufbereitung} im Anschluss an Marchi \autocite*{marchiDividingDataEpistemological2018} bereits angemerkt, bestehen zur Gliederung des Untersuchungskorpus in diachrone Segmente grundsätzlich verschiedene Herangehensweisen. Vielversprechend wäre es, Möglichkeiten der \emph{corpus-driven} Unterteilung anhand von sprachoberflächlichen, quantitativ messbaren Diskursveränderungen zu eruieren. So könnten zum einen neue Erkenntnisse zu diskursiven Umbrüchen gewonnen werden und zum anderen ein methodisches Problem gelöst werden, das die hier vorgenommene Unterteilung entlang fester Zeitabschnitte birgt. So fiel im Rahmen der Analyse mehrfach auf, dass diskursiv relevante Ereignisse, die an den Rändern der definierten Zeitabschnitte stattfanden, wie etwa die Publikation von ‚Deutschland schafft sich ab` oder der Anschlag vom Berliner Breitscheidplatz, in einzelnen Frames nicht deutlich emergent wurden, was auf die kleine zugrundeliegende Textmenge, in denen sie thematisiert wurden, zurückgeht.

Bezüglich der Gliederung des Korpus ist noch eine weitere Limitation der vorliegenden Arbeit zu benennen. So wurde das Korpus in den einzelnen Subkorpora entweder in Hinblick auf die diachrone oder die mediale Achse aggregiert. Das Subkorpus 1999--2001 umfasste also z.\,B. Texte aller drei Zeitungen, während sich das \textsc{Taz}-Subkorpus aus allen Zeiträumen speiste. Die Gliederung wurde so vorgenommen, um einerseits pro Subkorpus eine möglichst große Menge an Texten einbeziehen zu können, was beim Training der Word Embeddings zuträglich war, und andererseits aus forschungspraktischen Gründen, um nicht zu viele verschiedene Subkorpora zu bilden. Dadurch werden mögliche Unterschiede zwischen den aggregierten Akteuren bzw. Zeiträumen jedoch nivelliert und es wäre in einer folgenden Untersuchung von Interesse, z.\,B. einzelne Diskursereignisse -- sofern sie nicht in den Zeitungskorpora emergent wurden -- noch einmal im Kontrast der einzelnen Zeitungen zu analysieren.

Letztlich sei auf weitere methodische Erweiterungsmöglichkeiten verwiesen, die sich durch die Modellierung als Netzwerkgraph eröffnen. Hier könnten weitere algorithmische Möglichkeiten der Informationsanreicherung z.\,B. mit anderen \emph{community-detection}-Algorithmen oder Zentralitätsmaßen sowie deren graphischer Realisierung erprobt werden. Erste Experimente zur \emph{corpus-driven} Berechnung von besonders zentralen FE in den Frames fielen durchaus vielversprechend aus. So zeigte sich, dass in fast allen Subkorpora die Präposition \emph{gegen} ein eigenes, hochgradig im Graphen vernetztes FE (u.\,a. \cluster{99\_gegen}, \cluster{08\_gegen}, \cluster{20\_gegen}) darstellt, das Verbindungen in alle thematischen Bereiche aufweist (so werden Prozesse \emph{gegen} Angeklagte geführt, Vorwürfe des Antisemitismus \emph{gegen} Personen erhoben, der Kampf \emph{gegen} extremistische Bestrebungen aufgenommen, \emph{gegen} etwas demonstriert oder auch Krieg \emph{gegen} andere Länder geführt). Die zentrale Position von \emph{gegen} in den Frames ist dabei in Einklang mit Laschs \autocite*{laschZurVereinbarkeitDiskurslinguistisch2014} und Kalwas \autocite*{kalwaKonzeptIslamDiskurslinguistische2013} (korpuslinguistischen) Analysen, die ebenfalls auf die Rolle der Präposition im Kontext des Islam(ismus)diskurses hingewiesen haben und hebt den Wert datengeleiteter Verfahren bei der Identifikation von synsemantischen bzw. generell unauffälligen Sprachgebrauchsmustern hervor. Schließlich ist es -- auch in Anbetracht des sozialwissenschaftlich-definitorischen Gehaltes des Extremismuskonzeptes, der sich \emph{ex negativo} aus der Ablehnung der FdGO speist, -- durchaus passend, das Dagegengerichtetsein als Kernbestandteil der Extremismussemantik aufzufassen.

\section{Fazit}\label{fazit-8}

Die vorliegende Arbeit ist der Frage nachgegangen, ob und wie sprachoberflächliche, distributionelle Muster als semantische Frames lesbar gemacht werden können (Forschungsziel 1). Dabei ist es gelungen, wesentliche Frame-semantische Strukturen des Extremismusdiskurses anhand sprachoberflächlicher Distributions- bzw. Sprachgebrauchsmuster nachzuzeichnen. Dieser Nexus zwischen Form und Bedeutung ist in beiden Theorietraditionen angelegt. Aus Frame-semantischer Perspektive konstatiert Busse: „sprachlichen und textuellen Strukturen entsprechen Strukturen im (gesellschaftlichen) Wissen`` \autocite*[164]{busseDiskursSpracheGesellschaftliches2013} und aus korpuslinguistischer Sicht formuliert Tognini-Bonelli, dass „the choice patterns of words in text can create new, large and complex units of meaning`` \autocite*[101]{tognini-bonelliCorpusLinguisticsWork2001}. Meine Arbeit stellt den Versuch dar, diese Annahmen systematisch und synergetisch in die Praxis umzusetzen.

Dazu wurde eine \emph{corpus-driven} Methode entwickelt, die ohne die Notwendigkeit von Annotationen in der Lage ist, semantische Frames der Diskursebene bzw. der Ebene des gesellschaftlichen Wissens als komplexe, rekursive Gefüge von Frame-Elementen und Frame-Relationen zu erschließen und zu visualisieren (Forschungsziel 2). Unter einer korpuspragmatischen Perspektive bestand das Ziel dabei gerade nicht darin, kontextabstrakte Frames für \emph{das} Deutsche oder \emph{den} medialen Diskurs zu modellieren, sondern einzelne diachrone und akteurabhängige Framings zu kontrastieren.

Erprobungsgegenstand des Ansatzes und des entwickelten Verfahrens war das Framing einzelner Extremismusvarianten durch die Zeitungen \textsc{Taz}, \textsc{Spiegel} und \textsc{Welt} im Zeitraum 1999--2021 (Forschungsziel 3). Datengeleitet kristallisierten sich zunächst thematische Frame-Anteile heraus, die den Extremismusdiskurs als mit politischen, gesellschaftlichen, kriegerischen und juristischen Kontexten verwoben konturieren. Es konnten sodann einzelne Extremismusvarianten auf einer Mesoebene herausgearbeitet werden. Einerseits wurden so Konstanten des Framings einzelner Extremismen wie der Stigmacharakter des Extremismuskonzeptes, der sie als etwas zu Bekämpfendes semantisiert, oder die Organisiertheit in Gruppen sichtbar. Andererseits wurden im Kontrast auch Spezifika einzelner Formen deutlich, die sich etwa in Bezug auf die Neigung zu Gewalt und Terror zeigten. Im Fall des Islamismus ist diese konstant präsent; aus dem Framing von Rechtsextremismus hingegen verschwindet sie über die erste Dekade des 20. Jahrhunderts weitestgehend, kehrt dann aber deutlich zurück und im Falle des Linksextremismus, der -- außer in der \textsc{Taz} -- eher undeutlich in den Frames zu Tage tritt, wird Gewalt punktuell emergent und dann mit der RAF assoziiert oder als \emph{Krawalle} im Rahmen von Demonstrationen bezeichnet.

Insgesamt konnten mit der hier entwickelten Methode wesentliche Bedeutungsstrukturen eines aktuellen gesellschaftlichen Diskurses anhand distributioneller, \emph{corpus-driven} Verfahren empirisch nachgezeichnet und für eine Frame-semantische Analyse fruchtbar gemacht werden. Auch vor dem Hintergrund der jüngeren überwältigenden Entwicklungen im Bereich der \emph{large language models} wird so ein Beitrag geleistet, verständlich zu machen, warum und inwiefern diese Modelle, die doch nur auf der Grundlage von bloßen Textdaten und Statistik operieren, so viel über die Welt zu wissen scheinen.